% 第 2 章：交流序网法与故障类型
\section{交流序网法与故障类型}\label{sec:sc-ac}

\subsection{序阻抗矩阵}
IEC 60909 以故障处等效电压源替代故障前网络；非旋转负荷除非带显式电机短路数据否则忽略，
旋转机与构网源以内阻抗表示。对故障母线 $k$，构建序导纳矩阵 $Y_1,Y_2,Y_0$ 并求
序母线阻抗对角
\begin{equation}
Z_{1k}=Z_1(k,k),\quad Z_{2k}=Z_2(k,k),\quad Z_{0k}=Z_0(k,k),\qquad Z_s=Y_s^{-1}.
\end{equation}
故障阻抗在当前选项中为实值标幺电阻 $Z_f=R_f+j0$（\fld{fault\_impedance\_pu}）。故障类型由
\fld{FaultType}（\srcpath{short\_circuit.hpp:22}）选择。

\subsection{三相故障}
仅用正序网：
\begin{equation}
I_{k3}''=\frac{c}{Z_{1k}+Z_f}I_b,\qquad Z_{k,3}=Z_{1k}+Z_f,\qquad |I_{k3}''|=\frac{c}{|Z_{k,3}|}I_b.
\end{equation}
三相短路功率 $S_k''=\sqrt3\,U_b\,|I_{k3}''|$（\fld{BusFaultResult.sk\_mva}，:49）。

\subsection{单相接地故障}
序网串联：
\begin{equation}
I_{k1}''=\frac{3c}{Z_{1k}+Z_{2k}+Z_{0k}+3Z_f}I_b,\qquad
Z_{k,1}=\frac{Z_{1k}+Z_{2k}+Z_{0k}+3Z_f}{3}.
\end{equation}
当零序路径低阻时 $|(Z_1{+}Z_2{+}Z_0)/3|<|Z_1|$，故 $|I_{k1}''|>|I_{k3}''|$——有效接地系统
（强中性点或接地变路径）中单相电流可高于三相，非数值缺陷。

\subsection{两相故障}
\begin{equation}
I_{k2}''=\frac{\sqrt3\,c}{Z_{1k}+Z_{2k}}I_b,\qquad Z_{k,2}=\frac{Z_{1k}+Z_{2k}}{\sqrt3}.
\end{equation}

\subsection{两相接地故障}
实现用正序等效的紧致工程近似：
\begin{equation}
Z_{k,2E}=Z_{1k}+\big(Z_{2k}\parallel Z_{0k}\big)+Z_f,\quad
Z_{2k}\parallel Z_{0k}=\frac{Z_{2k}Z_{0k}}{Z_{2k}+Z_{0k}},\quad
|I_{k2E}''|=\frac{c}{|Z_{k,2E}|}I_b.
\end{equation}
若需相分解电流 $I_{L2}/I_{L3}$ 与地电流，须扩展为完整序电流重构（当前只给有效电流幅值）。

\subsection{概览路径与近似}
概览 \srcpath{compute\_short\_circuit}（:212）用 Eigen 复数 SparseLU 分解序导纳矩阵，对不平衡
故障采用简化 $Z_1{=}Z_2{=}Z_0$（负序=正序、零序未建模）；详细路径（第~\ref{sec:sc-detailed}~章）
构建真正/负/零序矩阵。
